\documentclass{article}
\usepackage[utf8]{inputenc}
\usepackage[T1]{fontenc}
\usepackage{geometry}
\usepackage{graphicx}
\usepackage{float}
\usepackage{hyperref}
\usepackage[english,indonesian]{babel}

\title{Optimasi Pemrosesan Template LaTeX Menggunakan Pendekatan Berbasis AI}
\author{Academic Support Agent}
\date{\today}

\begin{document}
\maketitle

\begin{abstract}

\end{abstract}

\chapter{PENDAHULUAN}


\chapter{TINJAUAN PUSTAKA}


% Original content (commented out):
% \documentclass[a4paper, 12pt, oneside]{article}
\usepackage{fzp_thesis}
\usepackage[hidelinks]{hyperref}
\usepackage{luatexja-fontspec}
\usepackage{calc}
\usepackage{pdfpages}
\usepackage{siunitx}
\usepackage{luavlna}
\sisetup{
    list-final-separator = { a },
    list-pair-separator = { a },
    range-phrase = {--},
    range-units = single,
    output-decimal-marker = {,},
}

\newcommand{\thesistitle}{Dlouhý název super práce o velmi zajímavých věcech}

\hypersetup{
    colorlinks=false,
    linkcolor=blue,
    filecolor=magenta,
    urlcolor=cyan,
    pdftitle={\thesistitle},
}

\setcounter{secnumdepth}{4}
\setcounter{tocdepth}{3}

\begin{document}

\pagenumbering{gobble}

\singlechars{czech}{AaIiVvOoUuSsZzKk}
\preventsinglelang{czech}
%\preventsingledebugon

\noindent
\includegraphics[width = 10cm]{figures/FZP_CZ_CZU.png}

\vspace{1.5cm}

\noindent
{\sffamily\fontsize{30}{30}\selectfont Dlouhý název super práce \\
o velmi zajímavých věcech\par}

\vspace{1.5cm}

\noindent
\textcolor{czu_gray_text}{\sffamily\large Jméno Příjmení}

\vspace{1.5cm}

\noindent
\textbf{\sffamily\large Bakalářská práce}

\vspace{5cm}

\begin{flushright}
\sffamily\textbf{Katedra ekologie\\
Fakulta životního prostředí\\
Česká zemědělská univerzita v Praze}
\vfill
Vedoucí práce: prof. Ing. Jméno Příjmení, CSc.

\vspace{\baselineskip}

Datum: 31.~3.~2023
\end{flushright}
\newpage

\includepdf[scale=0.9, pages=-]{zadani.pdf}

\section*{Čestné prohlášení}
Prohlašuji, že jsem bakalářskou práci na téma: \enquote{\thesistitle} vypracovala samostatně a citovala jsem všechny informační zdroje, které jsem v práci použila a které jsem rovněž uvedla na konci práce v seznamu použitých informačních zdrojů.

Jsem si vědoma, že na moji bakalářskou práci se plně vztahuje zákon č.~121/2000~Sb., o právu autorském, o právech souvisejících s právem autorským a o změně některých zákonů, ve znění pozdějších předpisů, především ustanovení §~35 odst. 3 tohoto zákona, tj. o užití tohoto díla.

Jsem si vědoma, že odevzdáním bakalářské práce souhlasím s jejím zveřejněním podle zákona č.~111/1998~Sb., o vysokých školách a o změně a doplnění dalších zákonů, ve znění pozdějších předpisů, a to i bez ohledu na výsledek její obhajoby.

Svým podpisem rovněž prohlašuji, že elektronická verze práce je totožná s verzi tištěnou a že 
s údaji uvedenými v práci bylo nakládáno v souvislosti s GDPR.

\vspace{2cm}
V Praze dne 31.~3.~2023
\newpage

\mbox{}
\vfill
{\normalfont\sffamily\Large\bfseries\color{fzp_main}\selectfont Poděkování}

Ráda bych poděkovala vedoucímu práce prof. Ing. Jméno Příjmení, CSc. 
Dále děkuji všem a hlavně sobě...


\section*{Abstrakt}
Bakalářská práce se zabývá spoustou zajímavých věcí.

Atd.

\noindent\textbf{Klíčová slova:} bakterie, viry, příroda

\section*{Abstract}
Abstract in english. Blah Blah Blah.

Something something.

\noindent\textbf{Keywords:} bacteria, virus, nature

\tableofcontents

\section{Úvod}
\pagenumbering{arabic}
 Lorem ipsum dolor sit amet, consectetuer adipiscing elit. Etiam quis quam. Suspendisse nisl. Etiam egestas wisi a erat. Pellentesque arcu. Integer tempor. Vivamus luctus egestas leo. Morbi leo mi, nonummy eget tristique non, rhoncus non leo. Donec iaculis gravida nulla. In dapibus augue non sapien.
 
 Etiam dictum tincidunt diam. Sed ut perspiciatis unde omnis iste natus error sit voluptatem accusantium doloremque laudantium, totam rem aperiam, eaque ipsa quae ab illo inventore veritatis et quasi architecto beatae vitae dicta sunt explicabo. Nemo enim ipsam voluptatem quia voluptas sit aspernatur aut odit aut fugit, sed quia consequuntur magni dolores eos qui ratione voluptatem sequi nesciunt. Fusce nibh. Sed ac dolor sit amet purus malesuada congue.

\section{Příklad sekce}
\subsection{Příklad podsekce s textem, jednotkami a citacemi}
Antarktický kontinent je oblast, do které spadá jak samotný kontinent, tak šelfové ledovce i Jižní oceán. Jedná se o čtvrtý největší kontinent světa o rozloze 13,830 milionů \unit{\km\squared}. Rozměry se ve zdrojích liší, ale publikace \parencite{prosek2013antarktida} se přiklání k údajům z British Antarctic Survey. Pro představu a porovnání je Evropa menší, její rozloha je 10,53 miliónů \unit{\km\squared}. Nezaledněná část Antarktidy zaobírá pouze \qty{44890}{\km\squared}. V moři se nachází v odhadu asi \qty{40}{\km} široký pruh, který je přirozenou hranicí této popisované oblasti. Nazýváme ho antarktická konvergence, nepřesně antarktická polární fronta. Nachází se mezi \ang{48} a \ang{61} jižní šířky. Antarktická konvergence sousedí s vodami Tichého, Atlantského a Indického oceánu. Sousedící oceány mají vodu teplejší a slanější, stýkají se s vodou poměrně sladkou a studenou. Změna teploty mezi oceány se pohybuje mezi \qtyrange{1}{5}{\celsius}. \parencite{prosek2013antarktida,peck2014oceanography}

\subsubsection{Příklad podpodsekce}
Nejvyšší vrchol kontinentu je Mt. Vinson. 4892~m~n.~m. Dalším důležitým místem jsou jižní póly. Jižní pól se nachází ve výšce 2835~m~n.~m. Jižní magnetický pól nemá stálou polohu. Jde o místo na jižní polokouli, kde jsou indukční čáry kolmé k Zemi. V průběhu roku se magnetický pól posouvá. Průměr posunu je \qty{12.5}{km} k severo-severozápadu. Jižní geomagnetický pól se také posouvá v průběhu roku v průměru o \qtyrange{2}{3}{\km} k jiho-jihozápadu. Oproti magnetickému pólu si lze geomagnetický pól představit jako osu s pomocí jednoduchého dipólu umístěném v geometrickém středu Země. Osa je odkloněna od osy rotace a protíná se se severním geomagnetickým pólem. \parencite{mcelhinny1998magnetic,prosek2013antarktida,wiki:xxx}

Hodnota se složitou jednotkou \qty{10}{\micro\mole\per\metre\squared\per\second}.

\subsection{Příklad dvou obrázků vedle sebe}
Dalším příkladem je obrázek \ref{fig:sklicko_pipetka}. Taky je tady text co zmiňuje pravý obrázek (obrázek \ref{fig:mikroskopovaci_stanice_2}).
\begin{figure}[h]
    \begin{minipage}[t]{.48\linewidth}
      \centering
      \includegraphics[width=\linewidth]{figures/R0003059.jpg}
      \caption{Podložní, krycí sklíčko, pipetka}
      \label{fig:sklicko_pipetka}
    \end{minipage}
    \hfill
    \begin{minipage}[t]{.48\linewidth}
      \centering
      \includegraphics[width=\linewidth]{figures/R0003057.jpg}
      \caption{Mikroskopovací stanice}
      \label{fig:mikroskopovaci_stanice_2}
    \end{minipage}
\end{figure}

Potom pokračuje text.


\section{Příklad jednoho obrázku}
Příklad reference na obrázek \ref{fig:schema1}.

\begin{figure}[h!]
      \centering
      \includegraphics[width=\linewidth]{figures/schema.pdf}
      \caption{Postup izolace na agarových plotnách}
      \label{fig:schema1}
\end{figure}

\section{Závěr}
 Nulla quis diam. Phasellus enim erat, vestibulum vel, aliquam a, posuere eu, velit. Fusce suscipit libero eget elit. Etiam posuere lacus quis dolor. Nullam rhoncus aliquam metus. Etiam neque. Cum sociis natoque penatibus et magnis dis parturient montes, nascetur ridiculus mus.
 
 Itaque earum rerum hic tenetur a sapiente delectus, ut aut reiciendis voluptatibus maiores alias consequatur aut perferendis doloribus asperiores repellat. Morbi scelerisque luctus velit. Nulla turpis magna, cursus sit amet, suscipit a, interdum id, felis. Sed elit dui, pellentesque a, faucibus vel, interdum nec, diam. Nullam sapien sem, ornare ac, nonummy non, lobortis a enim. Sed vel lectus. Donec odio tempus molestie, porttitor ut, iaculis quis, sem.


\newpage
\printbibliography[keyword={odborne},title={Odborné publikace}]
\printbibliography[keyword={internet},title={Internetové zdroje}]
\printbibliography[keyword={ostatni},title={Ostatní zdroje}]

\listoffigures
%\listoftables

\end{document}

\end{document}